%%%%%%%%%% TEMPLATE COREMA 2023%%%%%%%%%% %%%%%%%%%%%%%%%%%%%%%%%%%%%%%%%%%%%%%%%%
\documentclass[12pt, a4paper]{article}
\usepackage[left=3cm, top=3cm, right=2cm, bottom=2cm]{geometry}
\usepackage[utf8]{inputenc}
\usepackage[brazilian]{babel}
\usepackage{times}
\usepackage{color}
\usepackage{amsmath}
\usepackage{amssymb}
\usepackage{amsfonts}
\usepackage{mathrsfs}
\usepackage{graphicx}
\usepackage{textcomp, gensymb}
\usepackage{hyperref}
\usepackage{tikz}
\usepackage{emptypage}
\pagestyle{empty}

%\author{Comissão Científica COREMA 2023}
\begin{document}
%%%%%%%%%%%%%%%%%%%%%%%%%%%%%%%%%%
%%%% Não alterar, cabeçalho do Template

\begin{tikzpicture}[remember picture, overlay]
\node(Logo) at (current page.north west) [anchor=north west, xshift=3cm,yshift=-0.7cm]{\includegraphics[scale=0.05]{LOGO_OFICIAL_COREMA.png}};
\node at  (12,1) {Dias 04, 05 e 06 de setembro de 2023};
\node at  (10.35,0.4) {Campus de Engenharias e Ciências Agrárias -- Rio Largo};
\draw[black] (0,0)--(15,0);
\end{tikzpicture}
%%%%%%%%%%%%%%%%%%%%%%%%%%%%%%%%%%
    \begin{center}
    
    % Título
    \textbf{INSTRUÇÕES PARA ELABORAÇÃO DE TRABALHO PARA O CONGRESSO DE RECURSOS ENERGÉTICOS E MEIO AMBIENTE 2023 - Retificado3}
    \vspace{1cm}
    
    % Nome dos autores
    \textbf{Nome do primeiro autor$^1$ e Nome do segundo autor$^2$}\\
    $^1$Filiação, Cidade, E-mail do autor\\
    $^2$Filiação, Cidade, E-mail do autor
    \vspace{1cm}
    
    % Resumo
    \end{center}
    {\fontsize{10pt}{\baselineskip}\selectfont \textit{\textbf{\small{Resumo}}}}\\
    {\fontsize{10pt}{\baselineskip} \selectfont \textit{\small{Este documento serve como modelo para a preparação de trabalhos para o Congresso de Recursos Energéticos e Meio Ambiente 2023, apresentando instruções sobre as normas exigidas para submissões de artigos nos Anais do COREMA 2023. Siga as seguintes orientações: a) use no máximo 10 páginas e mínimo de 5, sendo estas no tamanho A4 (21 x 29,7 cm), com margens de 3 cm superior e a esquerda e 2 cm inferior e a direita; b) o corpo do texto deve ser digitado numa única coluna; c) não numerar as páginas ou colocar molduras; d) usar a fonte Times New Roman com tamanho 12 pt em todo corpo do texto; e) no resumo utilize no máximo 15 linhas, com fonte em itálico, 10 pt; f) use o alinhamento justificado (exceto no título e na descrição dos autores) e o espaçamento simples; g) as referências devem ser citadas de acordo com o título “Referências”, que está exibida no final deste documento; h) todos os símbolos e equações devem ser definidos; i) utilizar as unidades do Sistema Internacional de Unidades (SI); j) o resumo deverá conter no mínimo 2 palavras-chave; O trabalho deve ser escrito em português. O trabalho deverá ser submetido até o dia \textcolor{red}{06 de agosto}, devendo ser enviado no formato \textit{pdf.}, compatível com adobe acrobat e com o tema e subtema, conforme edital de submissão. O autor(a) deverá se inscrever no site www.coremaufal.com.br, para assim receber informações sobre sua submissão. O trabalho deve ser enviado através do site, na área de membros, na área de nome: ``minhas submissões".}}}
    \vspace{0.5cm}
    
    % Palavras-chave
    \noindent
    \textit{\textbf{\small{Palavras-chave: }}\small{Eficiência Energética, Aproveitamento de Resíduos, Reflorestamento (colocar no máximo 3 palavras-chave, separadas por vírgulas)}}
    \vspace{1cm}
    
% Introdução
\noindent
\textbf{1. \quad Introdução}
\vspace{0.5cm}

O trabalho deve seguir a seguinte estrutura de tópicos: 1. Introdução, 2. Fundamentação teórica, 3. Resultados e discussões, 4. Conclusões, e, por último, as Referências.

 Os trabalhos serão julgados por um Comitê Científico formada por professores e especialistas, sendo a relação de trabalhos aprovados divulgada após a avaliação por eles. Sendo o trabalho submetido, será analisado pelo Comitê Científico, que poderá aprovar ou solicitar alterações e/ou correções a serem feitas.
\vspace{0.5cm}

% Instruções de formatação
\noindent
\textbf{2. \quad Fundamentação teórica}
\vspace{0.5cm}

Na fundamentação teórica o autor deve inserir as teorias e modelos que fundamentam seu trabalho, discutindo os tópicos dos autores citados, revisando textos artigos, livros e outros materiais.
\vspace{0.5cm}

\noindent
2.1 \quad Seções e subseções
\vspace{0.5cm}

O título das seções deve ser digitado de tal forma que somente a primeira palavra tenha por sua primeira letra como maiúscula. Inicie o título com algarismos arábicos, seguido de um ponto, dando em seguida uma tabulação (apertando na tecla TAB do teclado). Digite a seção em negrito, alinhado à esquerda e com 12 pt. 

As subseções devem ser digitadas em 12 pt, de estilo normal, sem negrito ou itálico.
\vspace{0.5cm}

\noindent
\textbf{3. \quad Resultados e discussões}
\vspace{0.5cm}

Nesta seção coloque os seus resultados e discuta-os, podendo ser inseridos tabelas e gráficos como recursos de explicação. Siga conforme orientado nas seguintes subseções:
\vspace{0.5cm}

% Equações
\noindent
3.1 \quad Equações
\vspace{0.5cm}

As equações devem estar no centro da página, enumeradas em sequência, com os algarismos de numeração alinhados à direita. Deixe também um espaço antes e depois de cada equação. No Microsoft Word as equações devem ser feitas no ambiente ``Equação”. Prossiga como o exemplo a seguir:
\begin{equation}\label{Teorema do Transporte de Reynolds}
    \frac{DM_{sis}}{DT} = \frac{\partial}{\partial t}\int_{VC} \rho\;d\forall + \int_{SC} \rho \left( \vec{V}\hat{n}\right)\; dA 
\end{equation}

Sempre que a equação for chamada em alguma parte do texto, deve ser chamada como no formato supracitado, (\ref{Teorema do Transporte de Reynolds}).

Os símbolos das equações devem estar em itálico. Devendo ser definidos seu significado e unidades na primeira vez que forem mencionados. Todos os dados inseridos nos trabalhos devem estar no Sistema Internacional de Unidades (SI). A separação entre a parte inteira e decimal de números fracionários deve ocorrer por meio do uso de vírgulas.
\vspace{0.5cm}

% Tabelas
\noindent
3.2 \quad Tabelas
\vspace{0.5cm}

Todas as tabelas devem ser citadas no texto, numeradas pela ordem de aparição. A tabela deve estar centrada no texto. Posicione o título e a fonte, respectivamente acima e abaixo da tabela, deixando um espaço de uma linha. Deixe um espaço entre o corpo do texto e o título da tabela. Todas as tabelas devem ser referenciadas.

Para o título e fonte da tabela, bem como os dados nela contido utilize fonte 10 pt. Para os dados na tabela, coloque as respectivas unidades. Prossiga como no exemplo abaixo:

\begin{center}
\footnotesize{Tab. 1: Condições climáticas de Maceió – AL (normal climatológica 1981-2010)}
\vspace{0.2cm}

\begin{tabular}{|c|c|}
    \hline PARÂMETRO & VALOR \\\hline
    Temperatura máxima média anual (\celsius) & 29,8\\\hline
    Temperatura mínima média anual (\celsius) & 20,9\\\hline
    Temperatura média anual (\celsius) & 25,1\\\hline
    Precipitação anual (mm) & 1867,4\\ \hline
\end{tabular}
\vspace{0.2cm}

\footnotesize{Fonte: INMET.}
\end{center}
\vspace{0.5cm}

% Figuras
\noindent
3.3 \quad Figuras
\vspace{0.5cm}

Faça referência às figuras no corpo do texto. Assim como nas tabelas, coloque a fonte da figura. No caso das figuras, para facilitar a visualização, sugere-se o uso de imagens e gráficos de boa qualidade. Figuras devem estar localizadas no centro do texto. Deixe o espaço de uma linha de espaço entre as legendas e a tabela, reserve também uma linha de espaço entre a legenda e o texto. Numere as figuras, na ordem que aparecem no texto. Caso faça uso de gráficos, insira o significado de seus eixos coordenados, assim como as suas unidades. Faça como o seguinte gráfico:
\begin{center}
    \footnotesize{Fig. 1: a legenda deve ser escrita em fonte 10 pt}
    \vspace{0.2cm}
    
    \includegraphics[scale = 0.4]{índice.png}
    \vspace{0.2cm}
    
    \footnotesize{Fonte: o autor}
\end{center}
\vspace{0.5cm}

% Agradecimentos
\noindent
\textbf{4. \quad Conclusão}
\vspace{0.5cm}

Neste tópico faça as considerações finais do assunto abordado, reforçando a ideia principal desenvolvida.
\vspace{0.5cm}

% Referências
\noindent
\textbf{Referências}
\vspace{0.5cm}

Este espaço é reservado para inserção de cada referência utilizada na construção do
trabalho, de quaisquer que sejam as fontes. Todas as referências devem ser citadas no texto.
\vspace{0.5cm}

\noindent
Citações
\vspace{0.5cm}

\noindent
\quad Todas as citações no texto devem, de acordo com o número de autores:
\begin{enumerate}
    \item Apenas um(a) autor(a): inserir o nome do(a) autor(a) e o ano de publicação;
    \item Dois autores: os nomes dos autores e o ano de publicação;
    \item Três ou mais autores: insira o primeiro autor seguido por “et al.” e o ano de publicação.
\end{enumerate}
    Se for inserir “et al.” coloque de forma simples, sem itálico, conforme modelo abaixo.
    Faça as citações de forma direta, por exemplo, “No modelo unidimensional, NUNES et al. (2003) propõe...”
\vspace{0.5cm}

\noindent
Lista de referências
\vspace{0.5cm}

\textbf{Referências} deve estar em negrito, alinhado à esquerda e com tamanho 12 pt. Deixe uma linha de espaço entre o nome “Referências” e a lista de referências. Justifique as referências, cada referência da lista deve estar alinhada com os títulos, não necessitando de um TAB para começar a digitação. Cada referência deve ser separada por um espaço 12 pt.

Abaixo, há a exemplificação de como proceder, com referências de variados meios.

\noindent
Referências de livros:
\vspace{0.5cm}

\noindent
Munson, B.R., et al., 2004. Fundamentos de mecânica dos fluidos, Blucher, São Paulo.
\vspace{0.5cm}

\noindent
Hibbeler, R.C., 2011. Estática: mecânica para engenharia, Pearson Prentice Hall, São Paulo.
\vspace{0.5cm}

\noindent
Referências de jornais de publicações:
\vspace{0.5cm}

\noindent
Behar, O., 2018. Solar thermal power plants – A review of configurations and performance comparison, Renewable and Sustainable Energy Reviews, 92, 608-627.
\vspace{0.5cm}

\noindent
Referências de dissertações e teses:
\vspace{0.5cm}

\noindent
Zaversky, F., 2010. Solar Thermal Power Plants – Modelling of Parabolic Trough Collector Concepts, Thesis, Institute for Thermal Turbomachinery and Machine Dynamics, Graz University of Technology, Graz (Austria).
\vspace{0.5cm}

\noindent
Referências da web:
\vspace{0.5cm}

\noindent
Para referências da internet, deve ser fornecido o URL do site entre \url{www.minhareferencia.com}. Também deve-se inserir a data em que foi acessado pela última vez. Qualquer outra informação adicional (DOI, autores, e outras e referências) também tem de ser fornecida. As referências da web podem ser incluídas na lista de referências.
\vspace{-2cm}
% \renewcommand{\refname}{ }
% \bibliographystyle{plain}   
% \bibliography{referencias} 
\end{document}